\section{The two prime sums in the eighth period-height entry}

% BEGIN I88P1
\paragraph{I88P1. The inert primes with zero Frobenius label.}
Fix $0<\epsilon<1/52$, an original family $\mathcal F$,
$x\ge2$ with $\mathcal F.x_1<x$, a member $i$ and a lift $1\le j\le m$.
Put $q=x^\epsilon$, $b=\exp q$, $a=xb$, $X=X''(n)=a/w''(n)$,
$N=N_m$, $\eta=L.\mathrm{etaSz}$, $S=L.\mathrm{Sc}$ and $B_0=fmbq$.
For any set of distinct primes $P$ below or equal to $X$,
\[
 \sum_{p\in P}2\log p\,|E_{C,j}(p)|
 \le 18(1+1/\epsilon)B_0 .
\]
Here are the numerical consequences of the size windows that we
use for both prime sums:
\[
 k(b'+\eta)\le mb,\quad S,b'\le b,\quad
 X\eta\le b,\quad N^2,4N,q\le\eta,\quad
 f\eta\le b,\quad 2x\le b .
\]
Indeed the denominator of the defining $\ell_*$ is at least $k$,
since $n\ge1$ and $S+t+q\ge1$. Its window gives
$b'+S+\eta<\ell_*\le mb/k\le b$.
The node bounds include $4N^3+2t+6\le\eta/m$ and
$ktq\le\eta/m$. The size inequalities for $\eta$ give both
$x\eta\le w''(n)$ and $\eta\,t^3fk^3\le b$.
Thus $X\eta\le b$. Also $2x\le x\eta\le w''(n)\le b'\le b$.
The latter bound on $w''$ follows from $w''=b'/(xn)$.

Put $W_j=\sum_\tau w_j(\tau)$. The $f$ weights obey
$0\le w_j(\tau)\le 3b$. For thick places the window has upper
end $b+\eta+1$, for thin places $w''+\eta$, and for top
places $b_t+b'/x+\eta$. The bottom and density size conditions
give $b_t\le(1+1/4)b'$: the relative error in the bottom
inequality is at most $x^{-2}/4$. The lower endpoints in all
three windows are nonnegative. In particular at every place
with first coordinate false, thick places are excluded by the
definition of the thick embeddings, so $w_j\le 3(b'+\eta)$.

The original digit constant is $3$, giving
$E_{C,j}(p)=3fN^2\log a/\log p+W_j/(p^2-1)\ge0$.
We have $|P|N^2\le X\eta\le b$ and $p\le X\le a$.
Also $1/(p^2-1)\le2/p^2$ and
$\sum_{n\ge2}n^{-2}\le\sum_{n\ge2}(1/(n-1)-1/n)=1$.
After multiplying by $2\log p$ the two sums are at most
$6fb\log a$ and $12fb\log a$, respectively. Finally
$\log a=\log x+q\le(1+1/\epsilon)q$ since
$\epsilon\log x\le\exp(\epsilon\log x)$. This proves the claim.
% END I88P1

% BEGIN I88P2
\paragraph{I88P2. The other correction at one prime.}
For $p\in P$ let $d=n_P$, $\gamma=p/(p^2-1)$ and
$\delta_p={\bf1}_{\mathrm{IsCs}(p)}$, where $\mathrm{IsCs}(p)$
requires an inert prime with nonzero label $h_P$ of odd additive
order. The native constants are $c_6=12$, $d_{\ell,1}=5$ and
$d_{\ell,2}=4$. The local estimate is
\[
 |E_{I,j}(p)|\log p
 \le 126 fmb\frac{\log p}{p}
     +9fbk\delta_p\frac{\log p}{p^3}
     +70f\eta q .
\]
We give the bounds in the order in which the terms occur in the
definition of $E_I$. Each selected subset of the $f$ signed
places has size at most $f$. Also $1\le d\le f$, since $d$ is the
order of the label or twice that order, and the label group has
size $f/2$. In the first term the orbit optimization has upper
bound $\gamma\sum_s (k-1)w_j(z_s)$: its definition subtracts
the supremum of a nonnegative feasible sum from this expression.
The zero assignment is feasible and the feasible sums are
bounded by summing the constraints over the finite orbit.
The places $z_s$ in a kind-A orbit start with first coordinate false
at a nonsignal prime, so that coordinate stays false throughout.
By I88P1 each $(k-1)w_j(z_s)\le3k(b'+\eta)\le3mb$.
Sum over the selected places, divide by the orbit length $d$,
and use $\gamma\le2/p$. Taking the maximum with zero preserves the
bound $6fmb/p$.

The top-place term is at most $12fmb/p$, using
$b+\eta+1+3S\le6b$. The term with coefficient $12$ and
denominator $p$ is at most $108fmb/p$, since each kind-BC
weight is at most $9b$ (I88P1 even gives $3b$).
These three bounds account for the coefficient $126$.
For the next three terms use
\[
 \log(4dNa)\le3q,\qquad \log(2a)\le2q,\qquad
 1/2\le\log p\le q .
\]
Indeed $4dN\le f\eta\le b$, $2a=2xb\le b^2$ and $p\le X\le b$.
Multiplying by $\log p$, the adjacent-rank term is bounded by
$fN^2(3q+q)\le4f\eta q$. The logarithmic term with coefficient
$5$ is at most $10f\eta q$. The squared base-$p$ logarithm with
coefficient $4$ is at most
$4f(2q)^2/(1/2)\le32f\eta q$. This gives $46f\eta q$.
The remaining correction $E_X$ is supported on $\delta_p=1$.
Its two contributions after multiplication by $\log p$
are at most $9fbk\log p/p^3$ and $12fN^2\log a\le24f\eta q$.
All the stated terms are nonnegative with the native maximum in the
first term, so these estimates also bound the absolute value.
% END I88P2

% BEGIN I88P3
\paragraph{I88P3. Summation on the two Frobenius branches.}
Write $M=2^e$ for the two-primary parameter of the frame. At a
prime with $\delta_p=1$ its arithmetic Frobenius has second
component $\iota(h_P)$: the defining Frobenius reading is independent
of the prime above $p$ and the inert class has this component.
This second component has the odd order of $h_P$. Its image in the
real two-power cyclotomic quotient, a $2$-group, is trivial.
The first component, complex conjugation, acts trivially on that
real quotient as well.
On a primitive $2^{e+2}$-th root Frobenius acts by the $p$-th power.
On passing to the quotient of units by $\{1,-1\}$, triviality
says exactly $p\equiv\pm1\pmod{4M}$. These primes are unramified
in this cyclotomic reading, as prescribed by the inert class.

For either sign write $p=4Mv\pm1$ with $v\ge1$. Then $p\ge3Mv$,
$M\ge2$, and for $y=p/M\ge1$ we have
$\log p=\log M+\log y\le2y\log M$ by $\log y\le y-1$ and
$\log M\ge1/2$. Thus
\[
 \frac{\log p}{p^3}\le
    \frac{\log M}{M^3}\frac1{4v^2}.
\]
There are at most two such integers per $v$ and
$\sum_{v\ge1}v^{-2}\le2$. Moreover $k\le x^5$ by the rank
window $k\le x^{9/2+10\epsilon+\epsilon/100}$ and $\epsilon<1/52$.
The frame relation $n\le oM$, $o\le x^{\epsilon/100}$, and
$n\ge x^{2+8\epsilon-\epsilon/100}$ give $M\ge x^2$.
Consequently $k\le M^{5/2}$, and $\log M\le2\sqrt M$ gives
\[
 k\sum_{p\in P}\delta_p\frac{\log p}{p^3}\le2 .
\]

For the remaining prime sum let $h=\lfloor b\rfloor$. For every
$p\in P$, $p\le h$ and $v_p(h!)\ge\lfloor h/p\rfloor\ge h/(2p)$.
Taking logarithms in prime factorization gives
$\sum_{p\in P}\log p/p\le 2\log(h!)/h\le2q\le3q$.
Finally $|P|\eta\le X\eta\le b$. Summing I88P2 therefore gives
\[
 \sum_{p\in P}|E_{I,j}(p)|\log p
 \le378B_0+18fb+70fbq\le424B_0 ,
\]
where $m\ge2$ and $q\ge1$ make the last two terms at most
$9B_0$ and $35B_0$. Hence the original weighted sum is at most
$848B_0$. Partition any admissible $P$ of signal or inert primes
$p\le X$ according to whether it is inert with $h_P=0$.
On that part $E_{CI}=E_C$; on its complement $E_{CI}=E_I$.
Adding I88P1 and this bound gives the original inert-prime estimate
of H87P3 with $C_I=18(1+1/\epsilon)+848$.
The same threshold $\max(2,\mathcal F.x_1+1)$ works for all the
members, lifts and subsets of primes.
% END I88P3
